% SPDX-License-Identifier: CC-BY-NC-ND-4.0
% Copyright (c) 2026 Ingolf Lohmann.
% Local arXiv submission candidate. Not submitted.
\documentclass[11pt]{article}

\usepackage[T1]{fontenc}
\usepackage[utf8]{inputenc}
\usepackage{lmodern}
\usepackage{microtype}
\usepackage[a4paper,margin=27mm]{geometry}
\usepackage{amsmath,amssymb,amsthm,mathtools}
\usepackage{booktabs,longtable,tabularx,array}
\usepackage{enumitem}
\usepackage{xcolor}
\usepackage{hyperref}
\usepackage{xurl}
\usepackage[nameinlink,noabbrev]{cleveref}

\hypersetup{
  colorlinks=true,
  linkcolor=blue!50!black,
  citecolor=blue!50!black,
  urlcolor=blue!50!black,
  pdftitle={Observer-Relative Retrocausality in Distributed Information Systems},
  pdfauthor={Ingolf Lohmann},
  pdfsubject={Formal existence result, executable witnesses, and evidence-bound effect control},
  pdfkeywords={QIK-VRT, retrocausality, distributed systems, provenance, proper time, delayed choice, effect acknowledgement}
}

\setlength{\parindent}{0pt}
\setlength{\parskip}{0.45em}
\setlist[itemize]{leftmargin=1.5em,itemsep=0.15em,topsep=0.25em}
\setlist[enumerate]{leftmargin=1.8em,itemsep=0.2em,topsep=0.25em}

\newtheorem{definition}{Definition}[section]
\newtheorem{theorem}{Theorem}[section]
\newtheorem{proposition}{Proposition}[section]
\newtheorem{corollary}{Corollary}[section]
\theoremstyle{remark}
\newtheorem{remark}{Remark}[section]

\newcommand{\host}{\prec_{\!H}}
\newcommand{\receive}{\prec_{\!O}^{R}}
\newcommand{\sourceorder}{\prec^{\theta}}
\newcommand{\qikvrt}{QIK--VRT}
\newcommand{\zenodo}{\url{https://doi.org/10.5281/zenodo.21888130}}
\newcommand{\authoritysnapshot}{\url{https://github.com/Goldkelch/qik-vrt/tree/12842e8df99553260774d53517522b2b5539c8a8}}
\newcommand{\mirrorsnapshot}{\url{https://github.com/ingolf-lohmann/qik-vrt/tree/c52b324914978dc5d6d80251260ce55f396909f7}}

\title{Observer-Relative Retrocausality in Distributed Information Systems:\\
\large A Formal Existence Result, Executable Witnesses, and Evidence-Bound Effect Control}
\author{Ingolf Lohmann\\
Independent Researcher, QIK--VRT}
\date{12 August 2026}

\begin{document}
\maketitle

\begin{abstract}
We introduce \emph{observer-relative retrocausality}, an operational relation
in a distributed event system. Let $\tau_O$ be a strictly increasing local
time parameter for an observer $O$, and let $\theta$ be an authenticated,
order-preserving marker of source events. The relation occurs when two
information-bearing records are received in increasing local time while their
bound source markers decrease: $\Delta\tau_O>0$ and $\Delta\theta<0$.
Equivalently, reception order is inverse to source-event order for at least one
record pair. A short existence theorem shows that this relation is compatible
with an acyclic host order and with every individual record being emitted
before it is received. A finite executable witness and a positive-latency
construction make the relation reproducible.

The paper separates this operational result from stronger propositions that it
does not entail: proper time does not run backwards, no record is received
before its own emission, no past record is overwritten, and no controllable
backward signalling channel is introduced. Delayed-choice and quantum-eraser
experiments are discussed as a structural connection: later context can make a
classification of earlier records available without changing the earlier local
record or enabling signalling to the past. We also describe an evidence-profile
proposal for the Effect Acknowledgement protocol that records source reference,
receiver-local receipt order, and provenance before a downstream effect is
released. The paper is a concise English synthesis of the QIK--VRT research
programme; the wider universe-level correspondence is explicitly identified as
the author's thesis rather than as a consequence of the finite theorem.
\end{abstract}

\noindent\textbf{Keywords:} QIK--VRT; observer-relative retrocausality;
distributed systems; proper time; provenance; information-reference direction;
delayed choice; quantum eraser; effect acknowledgement.

\section{Introduction}
\label{sec:introduction}

Information systems routinely receive valid records in an order different from
the order in which their source events occurred. The familiar cause is unequal
latency: an older record can take a longer path, be delayed by storage or
relaying, and arrive after a newer record. The usual observation is mundane,
but it exposes a useful distinction that is often hidden by speaking simply of
`time`: source-event order, path order, local reception order, and the
evidence available to a particular observer are distinct relations.

This paper formalizes the distinction in the terminology of Ingolf Lohmann's
QIK--VRT research programme. It calls an inverse relation between an
observer's increasing local time and the authenticated source-time references
of newly received information \emph{observer-relative retrocausality}. The
term is defined operationally here. It must not be read as silently importing
every stronger use of retrocausality in the foundations of physics.

The central relation is deliberately compact:
\begin{equation}
  \Delta\tau_O>0 \qquad\text{and}\qquad \Delta\theta<0.
  \label{eq:core}
\end{equation}
The first inequality says that the observer advances in its own local time.
The second says that the later-received record refers to an earlier,
authenticated source event than the previously received record. The
backward direction is therefore a direction of \emph{information reference
relative to that observer}; it is not a reversal of the transport of an
individual signal.

The contributions are as follows.
\begin{enumerate}
  \item We give an exact, provenance-sensitive definition of observer-relative
  retrocausality and prove a finite existence theorem for it.
  \item We give an executable witness and a constructive positive-latency
  realization, including a second observer with the opposite reception order.
  \item We state the physical and epistemic boundaries: the result does not
  establish backwards proper time, past modification, causal loops,
  superluminal propagation, or controllable messages into the past.
  \item We connect the record/evidence structure to delayed-choice and
  quantum-eraser experiments without claiming that those experiments uniquely
  select QIK--VRT.
  \item We specify a companion evidence-profile proposal for effect control in
  distributed computing.
\end{enumerate}

The QIK--VRT programme, the operational terminology, and the author's broader
correspondence thesis originate with Ingolf Lohmann. The present manuscript
condenses and clarifies historical materials retained in the public QIK--VRT
corpus \cite{LohmannRoundTrip2026,LohmannRelationalTime2026,
LohmannDecisionSufficiency2026,LohmannMandelbrot2026,LohmannCanonicalMemory2026}.

\section{Event model and definition}
\label{sec:model}

We use a distributed event structure
\begin{equation}
  H=(E,\host),
\end{equation}
where $E$ is a set of events and $\host$ is a strict partial order. The
relation $e\host f$ means that $e$ precedes $f$ in the declared host or
effect-relevant order. It is irreflexive and transitive, so it has no causal
cycle. This is structurally analogous to the partial order of events in a
distributed system described by Lamport \cite{Lamport1978}.

An observer $O$ is represented by a chain $E_O\subseteq E$ and a strictly
increasing local-time map
\begin{equation}
  \tau_O:E_O\longrightarrow\mathbb{R},\qquad
  e\host f\Longrightarrow\tau_O(e)<\tau_O(f)\quad(e,f\in E_O).
  \label{eq:localtime}
\end{equation}
For a relativistic realization, $\tau_O$ may be proper time along a
future-directed timelike worldline. For a computational realization, it is a
calibrated monotone local receipt parameter. The formal theorem needs only the
monotonicity in \cref{eq:localtime}.

Let $S\subseteq E$ be a declared source chain. Its source marker
\begin{equation}
  \theta:S\longrightarrow\mathbb{R}
\end{equation}
is strictly order-preserving. A record $r$ contains a payload $p(r)$, a
source event $\sigma(r)\in S$, and an integrity/provenance witness $w(r)$
that binds the record to $\sigma(r)$. We abbreviate
\begin{equation}
  \theta(r)=\theta(\sigma(r)).
\end{equation}
The reception event of $r$ at $O$ is $\rho_O(r)\in E_O$. Every admissible
record satisfies
\begin{equation}
  \sigma(r)\host\rho_O(r).
  \label{eq:forwarddelivery}
\end{equation}
Thus the model explicitly requires source-before-reception for every single
record.

The provenance binding matters. Arbitrary labels written on unrelated
messages do not establish a temporal relation. A valid source marker must be
bound to a real source event and must preserve the stated source order. The
record must also be informational rather than a duplicate with no evidentiary
consequence. One useful operational test is strict refinement of the set of
histories compatible with the observer's evidence: if $\mathcal{K}_{k-1}$ is
the compatible set before accepting a record and $\mathcal{K}_k$ the set after
verification, the record is information-bearing when
\begin{equation}
  \mathcal{K}_k\subsetneq\mathcal{K}_{k-1}.
  \label{eq:information}
\end{equation}

On accepted records there are at least two relevant orders:
\begin{align}
 r\receive r' &\Longleftrightarrow
   \tau_O(\rho_O(r))<\tau_O(\rho_O(r')),\\
 r\sourceorder r' &\Longleftrightarrow
   \theta(r)<\theta(r').
\end{align}

\begin{definition}[Observer-relative retrocausality]
\label{def:orr}
Observer-relative retrocausality occurs for $O$ when there are authentic,
information-bearing records $r_1,r_2$ such that
\begin{equation}
  r_1\receive r_2\qquad\text{and}\qquad r_2\sourceorder r_1.
  \label{eq:inverseorders}
\end{equation}
When numerical markers are present, the associated comparative
information-reference slope is
\begin{equation}
  v_I^O(r_1,r_2)=
  \frac{\theta(r_2)-\theta(r_1)}
       {\tau_O(\rho_O(r_2))-\tau_O(\rho_O(r_1))}<0.
  \label{eq:slope}
\end{equation}
\end{definition}

The relation is observer-relative because the reception order is local to an
observer. It is not arbitrary: the records, source witnesses, and local
receipts can be checked. Different receivers can legitimately instantiate
different reception orders over the same source records.

\section{Existence theorem and constructive witness}
\label{sec:theorem}

\begin{theorem}[Existence of negative information-reference direction]
\label{thm:existence}
Let $s^-\host s^+$ be two events in the source chain with
$\theta(s^-)<\theta(s^+)$. Let $r^-$ and $r^+$ be authentic,
information-bearing records bound respectively to $s^-$ and $s^+$. Suppose
that their receiver-local delivery order is inverted,
\begin{equation}
  \rho_O(r^+)\host\rho_O(r^-),
  \label{eq:overtaking}
\end{equation}
and that each record separately satisfies \cref{eq:forwarddelivery}. Then
observer-relative retrocausality occurs for $O$. More specifically,
\begin{equation}
  \Delta\tau_O=
  \tau_O(\rho_O(r^-))-\tau_O(\rho_O(r^+))>0
\end{equation}
and
\begin{equation}
  \Delta\theta=\theta(r^-)-\theta(r^+)<0,
\end{equation}
so $v_I^O(r^+,r^-)<0$.
\end{theorem}

\begin{proof}
Because the reception events lie on the observer chain and
\cref{eq:overtaking} holds, strict monotonicity of $\tau_O$ gives
$\Delta\tau_O>0$. Because the source marker is order-preserving on the
source chain, $\theta(s^-)<\theta(s^+)$, so $\Delta\theta<0$. The quotient
in \cref{eq:slope} is therefore negative. The individual forward-delivery
conditions continue to hold, so the inverse relation is between source order
and reception order, not within an individual signal path.
\end{proof}

The proof is intentionally short. It is a conditional theorem about an
explicit relation, not a claim that an arbitrary out-of-order packet proves a
new dynamical law. Its utility is to make the exact conditions inspectable:
source binding, information content, monotone local time, and inverse
reception/source order.

\begin{proposition}[No causal loop follows]
\label{prop:noloop}
The assumptions of \cref{thm:existence} do not imply a backward delivery of a
single record or a causal loop in $H$.
\end{proposition}

\begin{proof}
For every record, \cref{eq:forwarddelivery} gives
$\sigma(r)\host\rho_O(r)$. Backward delivery of that same record would also
require $\rho_O(r)\host\sigma(r)$, contradicting irreflexivity by
transitivity. Since $\host$ is a strict partial order, it has no closed
causal chain.
\end{proof}

\subsection{Finite witness}

Take two source events with markers $\theta(s^-)=10$ and
$\theta(s^+)=20$. Let their reception events occur at receiver-local times
$100$ and $101$ in the inverse order:
\begin{equation}
  \tau_O(\rho_O(r^+))=100,\qquad
  \tau_O(\rho_O(r^-))=101.
\end{equation}
Then
\begin{equation}
  v_I^O(r^+,r^-)=\frac{10-20}{101-100}=-10<0.
\end{equation}
The associated witness can be made information-bearing by beginning with four
possible two-bit histories. The first received record binds one bit and
reduces the compatible histories from four to two; the second binds the other
bit and reduces them from two to one. Equation \cref{eq:information} then
holds for both records.

The companion QIK--VRT corpus contains a network-independent checker,
\path{verify_observer_relative_retrocausality.py}, and its canonical
JSON witness report. It checks acyclicity, source-before-reception,
monotonicity, strict evidence refinement, inverse source/reception order, a
second observer, and a positive-latency realization. The exact source
snapshots are linked in \cref{sec:reproducibility}. The checker is an
executable finite witness, not a substitute for a proof assistant kernel or an
empirical survey of nature.

\subsection{Positive-latency realization}

Let the older record be emitted at coordinate time $t^-=0$ with delay
$d^-=3$, and the newer record at $t^+=1$ with delay $d^+=1$. Their reception
times are
\begin{equation}
  T^-=t^-+d^-=3,\qquad T^+=t^++d^+=2.
\end{equation}
Both delays are positive, yet the later source event is received first. In a
relativistic idealization with propagation at or below $c$, write
$d^\pm=L^\pm/c$. The condition for overtaking is
\begin{equation}
  \frac{L^- - L^+}{c}>t^+-t^-.
  \label{eq:latencycondition}
\end{equation}
This can be implemented by unequal optical paths, network routes, buffers, or
relays. It realizes the operational relation in \cref{def:orr}; it does not
make any path past-directed.

\section{Observer perspective and proper time}
\label{sec:observers}

The same source records can have different reception orders for different
observers. Receiver $O$ can receive $r^+$ before $r^-$ and satisfy
$v_I^O<0$, while $O'$ receives $r^-$ before $r^+$ and satisfies
$v_I^{O'}>0$. There is no contradiction because the two inequalities refer to
different local reception chains. What remains common is the authenticated
source order and the declared host order.

In a physical relativistic realization, proper time remains increasing along a
future-directed timelike worldline. For motion with $|v|<c$ in an inertial
frame,
\begin{equation}
  d\tau=dt\sqrt{1-v^2/c^2}>0.
\end{equation}
This is why the formal relation should not be described as proper time
running backwards. Rather, proper time supplies the fixed local direction
against which the source-reference direction of the received information is
compared. The distinction is compatible with the relativistic treatment of
local clocks and event order \cite{Einstein1905,Minkowski1909}.

\section{Delayed choice, quantum erasure, and later evidence}
\label{sec:quantum}

Delayed-choice and quantum-eraser experiments are relevant here as a
structural connection, not as a unique proof of the QIK--VRT interpretation.
In the delayed-choice quantum eraser of Kim et al., earlier signal detections
can later be sorted with suitable partner and context records into correlation
classes showing different conditional patterns \cite{ScullyDruehl1982,Kim2000}.
For complementary conditional patterns one may schematically write
\begin{align}
  P(x\mid j=1)&=\frac{1+\cos\phi(x)}{2},\\
  P(x\mid j=2)&=\frac{1-\cos\phi(x)}{2}.
\end{align}
At equal weights their unconditioned local marginal is
\begin{equation}
  P(x)=\tfrac12P(x\mid j=1)+\tfrac12P(x\mid j=2)=\tfrac12.
\end{equation}
Thus the earlier local record does not carry a selectable backward signal
\cite{Jordan1983}. The later context makes a correlation classification of
the earlier record available; it does not rewrite the earlier raw data.

Related delayed-choice realizations and delayed-choice entanglement swapping
make the same separation useful: later measurement context and joint data
analysis can change which conditional relation an observer is justified in
asserting about earlier registrations, without supplying a controllable
message into the observer's causal past \cite{Jacques2007,Ma2012,Ma2013}.
The experiments therefore exhibit a physical family of later evidence about
earlier records situations. They do not uniquely select QIK--VRT over other
interpretations, and they do not establish arbitrary physical retrocausality.
The historical QIK--VRT decision-sufficiency paper makes the same distinction
between later evidence classification and a retrospective mutation of a
physical record \cite{LohmannDecisionSufficiency2026}.

\section{Evidence-bound effect control}
\label{sec:eap}

The formal distinction between reception order and source reference has a
direct systems-engineering implication. A receiver should not treat a
received item as decision-sufficient solely because it arrived. Before an
irreversible or consequential effect is released, the evaluator may need to
know: What source event does this record represent? Is that binding
authenticated? Which local reception sequence applies? Is later evidence
required to classify the record correctly? Has the evidence chain been
preserved?

The companion proposal is named the \emph{EAP Temporal-Provenance and
Reference Profile} (EAP--TPRP). It is an external evidence profile for
\texttt{EFFECT\_ACK} version 1. It introduces no version-1 wire fields and no
change to the version-1 release predicate. It binds, as external evidence:
\begin{enumerate}
  \item an authenticated \texttt{source\_event\_id} and
  \texttt{source\_order\_marker};
  \item a \texttt{receiver\_id} plus
  \texttt{receiver\_receipt\_sequence} or \texttt{receipt\_timestamp};
  \item an immutable provenance commitment and evidence references; and
  \item one of \texttt{FORWARD\_REFERENCE},
  \texttt{RETROGRADE\_REFERENCE}, or \texttt{INDETERMINATE}.
\end{enumerate}
Here \texttt{RETROGRADE\_REFERENCE} means only that, for one receiver, a
later locally received authenticated evidence object refers to an earlier
source-order marker than a previously received object. It does \emph{not}
assert backward signalling, alteration of past events, physical or ontic
retrocausality, truth of payload content, or sender intent.

This is a proposed application-layer control and provenance layer around a
stored-program computer architecture. It is not a replacement for the
von-Neumann architecture, a new processor instruction, or an IETF-recognized
extension of a computing architecture. Its intended engineering benefit is
to make the evidence and origin conditions of a proposed effect auditable
before release. In distributed environments, that can help expose
information asymmetries; it cannot by itself prove exploitation, motive, or
truth.

At the time of this manuscript, the public base document
\path{draft-lohmann-qikvrt-effect-ack-03} is an active individual Internet
Draft, not an RFC, IETF standard, working-group product, consensus result, or
IETF endorsement \cite{LohmannEAP2026}. EAP--TPRP is a local,
not-yet-submitted companion candidate.

\section{Scope, falsifiability, and correspondence}
\label{sec:scope}

The operational result is falsifiable for any claimed instance. It fails if
any of the following conditions is absent:
\begin{itemize}
  \item the source markers are freely assigned, unauthenticated, or not
  order-preserving on the declared source chain;
  \item the two receipts are not on the same monotone observer chain;
  \item a purported record is not information-bearing in the sense of
  \cref{eq:information};
  \item the provenance binding can be silently altered; or
  \item the supposedly inverse relation disappears under the declared
  order-preserving reparameterizations.
\end{itemize}

The theory is therefore not the proposition that any delayed packet is a
time machine. It is a defined and checkable relation between two bound
orders. The result establishes a formal consequence of the stated model and
a constructive realization of that model's relation. A run on real hardware
is a physical instance of the computational/record relation if the declared
reference bindings are actually satisfied. Such a run is not, by itself, an
independent confirmation that all natural systems or the universe have the
same structure.

Lohmann's broader QIK--VRT correspondence thesis is that the universe itself
is fundamentally a distributed relational system and that this framework
describes reality in its stated model scope \cite{LohmannRoundTrip2026}. This
paper preserves that thesis as the author's explicitly attributed position. It
does not present the thesis as scientific consensus or as a deduction from the
finite theorem alone.

\section{Reproducibility and historical continuity}
\label{sec:reproducibility}

The complete historical QIK--VRT corpus, including earlier intermediate
documents and machine-readable materials, is publicly preserved at Zenodo
under \zenodo. The two retained historical intermediate papers document the
development of relational time and decision sufficiency; the preceding
Mandelbrot and canonical-memory works document related distinctions between
classification, provenance, and effect acknowledgement
\cite{LohmannRelationalTime2026,LohmannDecisionSufficiency2026,
LohmannMandelbrot2026,LohmannCanonicalMemory2026}. The immutable source
snapshots used for reproducibility are available from the Authority repository
at \authoritysnapshot{} and the Mirror repository at \mirrorsnapshot{}.
These references are immutable snapshots, not assertions that either
repository state presently represents a new publication or an IETF submission.

The source package for this manuscript contains only the concise English
synthesis. Its accompanying submission manifest records the artifact digests
and explicitly states that it is a local candidate until a separate
artifact-bound publication authorization and a platform submission are
completed. Earlier QIK--VRT PDFs are retained as historical intermediate
states rather than overwritten. This preserves the research history while
separating it from the current clarification.

\section{Conclusion}

In a distributed system an observer can advance monotonically in local time
while the authenticated source-event references of its newly received records
run in the opposite direction. Under the operational definition in this
paper, that is observer-relative retrocausality:
\begin{equation}
  \Delta\tau_O>0\quad\text{and}\quad\Delta\theta<0.
\end{equation}
The relation is mathematically transparent, constructively realizable with
positive delays, and testable with provenance-bound records. It neither
requires nor implies a reversed clock, a signal travelling into the past, or a
causal loop. Its scientific value is the precise separation of source order,
receiver-local order, and later evidentiary classification. Its engineering
value is a proposed way to bind those distinctions to effect-control decisions
without turning receipt alone into authorization.

\section*{Author contribution and disclosure}

Ingolf Lohmann originated the QIK--VRT research programme, the stated
operational definition, and the correspondence thesis. This manuscript is a
new English synthesis of that programme. Artificial-intelligence assistance
was used for drafting, formal organization, consistency checking, and layout;
the distinction between author-attributed claims, formal results, and open
scientific scope is retained explicitly.

\begin{thebibliography}{99}

\bibitem{Einstein1905}
A. Einstein, Zur Elektrodynamik bewegter Korper, \emph{Annalen der
Physik} 322, 891--921 (1905).
\href{https://doi.org/10.1002/andp.19053221004}{doi:10.1002/andp.19053221004}.

\bibitem{Minkowski1909}
H. Minkowski, Raum und Zeit, \emph{Physikalische Zeitschrift} 10,
104--111 (1909).

\bibitem{Lamport1978}
L. Lamport, Time, Clocks, and the Ordering of Events in a Distributed
System, \emph{Communications of the ACM} 21, 558--565 (1978).
\href{https://doi.org/10.1145/359545.359563}{doi:10.1145/359545.359563}.

\bibitem{ScullyDruehl1982}
M. O. Scully and K. Druhl, Quantum eraser: A proposed photon correlation
experiment concerning observation and delayed choice in quantum mechanics,
\emph{Physical Review A} 25, 2208--2213 (1982).
\href{https://doi.org/10.1103/PhysRevA.25.2208}{doi:10.1103/PhysRevA.25.2208}.

\bibitem{Kim2000}
Y.-H. Kim, R. Yu, S. P. Kulik, Y. H. Shih, and M. O. Scully, Delayed
Choice Quantum Eraser, \emph{Physical Review Letters} 84, 1--5 (2000).
\href{https://doi.org/10.1103/PhysRevLett.84.1}{doi:10.1103/PhysRevLett.84.1}.

\bibitem{Jacques2007}
V. Jacques et al., Experimental Realization of Wheeler's Delayed-Choice
Gedanken Experiment, \emph{Science} 315, 966--968 (2007).
\href{https://doi.org/10.1126/science.1136303}{doi:10.1126/science.1136303}.

\bibitem{Ma2012}
X.-S. Ma et al., Experimental delayed-choice entanglement swapping,
\emph{Nature Physics} 8, 479--484 (2012).
\href{https://doi.org/10.1038/nphys2294}{doi:10.1038/nphys2294}.

\bibitem{Ma2013}
X.-S. Ma et al., Quantum erasure with causally disconnected choice,
\emph{Proceedings of the National Academy of Sciences} 110, 1221--1226
(2013). \href{https://doi.org/10.1073/pnas.1213201110}{doi:10.1073/pnas.1213201110}.

\bibitem{Jordan1983}
T. F. Jordan, Quantum correlations do not transmit signals,
\emph{Physics Letters A} 94, 264 (1983).
\href{https://doi.org/10.1016/0375-9601(83)90713-2}{doi:10.1016/0375-9601(83)90713-2}.

\bibitem{LohmannRoundTrip2026}
I. Lohmann, Von Softwarearchitektur zur Weltformel -- DAS UNIVERSUM ALS
ROUND TRIP, version 1.0.0-candidate.1, Zenodo (2026).
\href{https://doi.org/10.5281/zenodo.21888130}{doi:10.5281/zenodo.21888130}.

\bibitem{LohmannRelationalTime2026}
I. Lohmann, Relationale Zeit, virtuelle Retrokausalitat und die monoton
wachsende Evidenzkugel: Ein beweisbares QIK--VRT-Modell mit massivem Kern,
unscharfem Rand und expliziter physikalischer Korrespondenzgrenze,
QIK--VRT historical intermediate manuscript (11 August 2026), preserved in
the associated QIK--VRT publication package \cite{LohmannRoundTrip2026}.

\bibitem{LohmannDecisionSufficiency2026}
I. Lohmann, Decision Sufficiency, Delayed Choice, and Witness-Bound Recovery:
QIK--VRT: formal boundaries, finite verification, and provenance-bound
continuation, QIK--VRT historical intermediate manuscript (6 August 2026),
preserved as a byte-exact historical intermediate state in the associated
QIK--VRT publication package.

\bibitem{LohmannMandelbrot2026}
I. Lohmann, Mandelbrot-Menge, Anschlussordnung und physikalisches
Modelluniversum -- Komplementbeweis, rekursive Wirkungsklassifikation,
dimensionsrichtige Raumzeitbruecke, Retrokausalitat und der entscheidende
Unterschied zur empirischen Quantengravitation, version 3.0, Zenodo (2026).
\href{https://doi.org/10.5281/zenodo.21482023}{doi:10.5281/zenodo.21482023}.

\bibitem{LohmannCanonicalMemory2026}
I. Lohmann, QIK--VRT und das Effect-Acknowledgement-Protokoll: Kanonischer
Speicher zwischen Vergangenheit und Zukunft, version 1.0.0, Zenodo.
\href{https://doi.org/10.5281/zenodo.21711193}{doi:10.5281/zenodo.21711193}.

\bibitem{LohmannEAP2026}
I. Lohmann, QIK--VRT Effect Acknowledgement: Separating Receipt from
Authorization for Downstream Effect, Internet-Draft
\texttt{draft-lohmann-qikvrt-effect-ack-03}, active individual submission
(2026). \url{https://datatracker.ietf.org/doc/draft-lohmann-qikvrt-effect-ack/}.

\end{thebibliography}

\end{document}
